% GENERATED FILE. Edit canonical lessons, then run npm run generate:books.
% canonical-source-hash: fnv1a64:f0c2d4ea8b799cc8
% canonical-lessons: MR-C307-prarambh, MR-C307-ant, MR-C307-akher, MR-C307-kalkhand, MR-C307-smrti

\chapter{Beginnings and Endings}
\label{ch:mr-beginnings-and-endings}
\hlchaptermodality{307}

\begin{chapteropening}
\textbf{By the end of this chapter:} \emph{I can say how a story or a period began and ended.}

Complete the last lesson of chapter 307: I can say how a story or a period began and ended.
\end{chapteropening}

\section[\texorpdfstring{\mr{प्रारंभ} (prārambh)}{प्रारंभ (prārambh)}]{\mr{प्रारंभ} (prārambh) — a commencement}

\label{lesson:MR-C307-prarambh}



\begin{quote}
Before the new one: say the Marathi for by itself, then the Marathi for lately.
\end{quote}

\subsection*{You'll want to know: \mr{प्रारंभ}}
\textbf{\mr{प्रारंभ}} — \emph{prārambh} — \textquotedblleft{}a commencement\textquotedblright{}.

\begin{quote}
\mr{कार्यक्रमाचा} \mr{प्रारंभ} \mr{गाण्याने} \mr{झाला}.
\end{quote}

\begin{quote}
\emph{kāryakramācā prārambh gāṇyāne jhālā.}
\end{quote}

\begin{quote}
The programme began with a song.
\end{quote}

Masculine.

The formal word for \mr{सुरुवात}.

\begin{cousinweb}
\textbf{\mr{प्रारंभ}} (\emph{prārambh}) is a tatsama, taken from Sanskrit \emph{prārambha}, 'a commencement, a beginning' (Monier-Williams).
\end{cousinweb}

\begin{grammarlens}[title={What you've built}]
One more word for telling what happened.
\end{grammarlens}

\subsection*{Your turn}
\begin{itemize}
\raggedright
  \item \emph{Say it:} \emph{prārambh}
  \item \emph{Say it:} \emph{prārambh}, once more
  \item \emph{Say it:} \emph{kāryakramācā prārambh gāṇyāne jhālā.} — The programme began with a song.
  \item \emph{Recall:} say the Marathi for by itself, then the Marathi for lately, then say \emph{prārambh} again
\end{itemize}

\begin{tcolorbox}[breakable,colback=teal!4,colframe=teal!35!black,title={Before you move on}]
What does \mr{प्रारंभ} mean? (\textquotedblleft{}A commencement\textquotedblright{}.) Say it once more.
\end{tcolorbox}
\section[\texorpdfstring{\mr{अंत} (ant)}{अंत (ant)}]{\mr{अंत} (ant) — an ending; an end}

\label{lesson:MR-C307-ant}



\begin{quote}
Before the new one: say the Marathi for lately, then the Marathi for a commencement.
\end{quote}

\subsection*{You'll want to know: \mr{अंत}}
\textbf{\mr{अंत}} — \emph{ant} — \textquotedblleft{}an ending; an end\textquotedblright{}.

\begin{quote}
\mr{या} \mr{कथेचा} \mr{अंत} \mr{फार} \mr{अनपेक्षित} \mr{होता}.
\end{quote}

\begin{quote}
\emph{yā kathecā ant phār anpekṣit hotā.}
\end{quote}

\begin{quote}
The ending of this story was very unexpected.
\end{quote}

Masculine.

\mr{शेवट} is the everyday word; \mr{अंत} is also the end of a life.

\begin{cousinweb}
\textbf{\mr{अंत}} (\emph{ant}) is a tatsama, taken from Sanskrit \emph{anta}, 'an end, a limit' (Monier-Williams).
\end{cousinweb}

\begin{grammarlens}[title={What you've built}]
One more word for telling what happened.
\end{grammarlens}

\subsection*{Your turn}
\begin{itemize}
\raggedright
  \item \emph{Say it:} \emph{ant}
  \item \emph{Say it:} \emph{ant}, once more
  \item \emph{Say it:} \emph{yā kathecā ant phār anpekṣit hotā.} — The ending of this story was very unexpected.
  \item \emph{Recall:} say the Marathi for lately, then the Marathi for a commencement, then say \emph{ant} again
\end{itemize}

\begin{tcolorbox}[breakable,colback=teal!4,colframe=teal!35!black,title={Before you move on}]
What does \mr{अंत} mean? (\textquotedblleft{}An ending; an end\textquotedblright{}.) Say it once more.
\end{tcolorbox}
\section[\texorpdfstring{\mr{अखेर} (akher)}{अखेर (akher)}]{\mr{अखेर} (akher) — at last; finally}

\label{lesson:MR-C307-akher}



\begin{quote}
Before the new one: say the Marathi for a commencement, then the Marathi for an ending.
\end{quote}

\subsection*{You'll want to know: \mr{अखेर}}
\textbf{\mr{अखेर}} — \emph{akher} — \textquotedblleft{}at last; finally\textquotedblright{}.

\begin{quote}
\mr{तीन} \mr{दिवसांनी} \mr{अखेर} \mr{पाऊस} \mr{थांबला}.
\end{quote}

\begin{quote}
\emph{tīn divsānnī akher pāūs thāmblā.}
\end{quote}

\begin{quote}
After three days the rain finally stopped.
\end{quote}

Also a noun: \mr{महिन्याच्या} \mr{अखेरीस}, at the end of the month.

\begin{cousinweb}
\textbf{\mr{अखेर}} (\emph{akher}) is a loan from Persian \emph{āxir}, 'the end', usually traced to Arabic \emph{ākhir}, 'last'.
\end{cousinweb}

\begin{grammarlens}[title={What you've built}]
One more word for telling what happened.
\end{grammarlens}

\subsection*{Your turn}
\begin{itemize}
\raggedright
  \item \emph{Say it:} \emph{akher}
  \item \emph{Say it:} \emph{akher}, once more
  \item \emph{Say it:} \emph{tīn divsānnī akher pāūs thāmblā.} — After three days the rain finally stopped.
  \item \emph{Recall:} say the Marathi for a commencement, then the Marathi for an ending, then say \emph{akher} again
\end{itemize}

\begin{tcolorbox}[breakable,colback=teal!4,colframe=teal!35!black,title={Before you move on}]
What does \mr{अखेर} mean? (\textquotedblleft{}At last; finally\textquotedblright{}.) Say it once more.
\end{tcolorbox}
\section[\texorpdfstring{\mr{कालखंड} (kālkhaṇḍ)}{कालखंड (kālkhaṇḍ)}]{\mr{कालखंड} (kālkhaṇḍ) — a period; an era}

\label{lesson:MR-C307-kalkhand}



\begin{quote}
Before the new one: say the Marathi for an ending, then the Marathi for at last.
\end{quote}

\subsection*{You'll want to know: \mr{कालखंड}}
\textbf{\mr{कालखंड}} — \emph{kālkhaṇḍ} — \textquotedblleft{}a period; an era\textquotedblright{}.

\begin{quote}
\mr{तो} \mr{कालखंड} \mr{देशासाठी} \mr{फार} \mr{कठीण} \mr{होता}.
\end{quote}

\begin{quote}
\emph{to kālkhaṇḍ deśāsāṭhī phār kaṭhīṇ hotā.}
\end{quote}

\begin{quote}
That period was very hard for the country.
\end{quote}

Masculine.

Built from \mr{काल}, time, and \mr{खंड}, a piece.

\begin{cousinweb}
\textbf{\mr{कालखंड}} (\emph{kālkhaṇḍ}) is a learned compound built from Sanskrit \emph{kāla}, 'time', and \emph{khaṇḍa}, 'a piece, a section' (Monier-Williams).
\end{cousinweb}

\begin{grammarlens}[title={What you've built}]
One more word for telling what happened.
\end{grammarlens}

\subsection*{Your turn}
\begin{itemize}
\raggedright
  \item \emph{Say it:} \emph{kālkhaṇḍ}
  \item \emph{Say it:} \emph{kālkhaṇḍ}, once more
  \item \emph{Say it:} \emph{to kālkhaṇḍ deśāsāṭhī phār kaṭhīṇ hotā.} — That period was very hard for the country.
  \item \emph{Recall:} say the Marathi for an ending, then the Marathi for at last, then say \emph{kālkhaṇḍ} again
\end{itemize}

\begin{tcolorbox}[breakable,colback=teal!4,colframe=teal!35!black,title={Before you move on}]
What does \mr{कालखंड} mean? (\textquotedblleft{}A period; an era\textquotedblright{}.) Say it once more.
\end{tcolorbox}
\section[\texorpdfstring{\mr{स्मृती} (smṛtī)}{स्मृती (smṛtī)}]{\mr{स्मृती} (smṛtī) — a remembrance; memory}

\label{lesson:MR-C307-smrti}



\begin{quote}
Before the new one: say the Marathi for at last, then the Marathi for a period.
\end{quote}

\subsection*{You'll want to know: \mr{स्मृती}}
\textbf{\mr{स्मृती}} — \emph{smṛtī} — \textquotedblleft{}a remembrance; memory\textquotedblright{}.

\begin{quote}
\mr{आजोबांच्या} \mr{स्मृती} \mr{अजून} \mr{ताज्या} \mr{आहेत}.
\end{quote}

\begin{quote}
\emph{ājobāñcyā smṛtī ajūn tājyā āhet.}
\end{quote}

\begin{quote}
Memories of Grandfather are still fresh.
\end{quote}

Feminine.

\mr{स्मृतीत}, in memory. \mr{आठवण} is the everyday word.

\begin{cousinweb}
\textbf{\mr{स्मृती}} (\emph{smṛtī}) is a tatsama, taken from Sanskrit \emph{smṛti}, 'remembrance, memory' (Monier-Williams).
\end{cousinweb}

\begin{grammarlens}[title={What you've built}]
One more word for telling what happened.
\end{grammarlens}

\subsection*{Your turn}
\begin{itemize}
\raggedright
  \item \emph{Say it:} \emph{smṛtī}
  \item \emph{Say it:} \emph{smṛtī}, once more
  \item \emph{Say it:} \emph{ājobāñcyā smṛtī ajūn tājyā āhet.} — Memories of Grandfather are still fresh.
  \item \emph{Recall:} say the Marathi for at last, then the Marathi for a period, then say \emph{smṛtī} again
\end{itemize}

\begin{tcolorbox}[breakable,colback=teal!4,colframe=teal!35!black,title={Before you move on}]
What does \mr{स्मृती} mean? (\textquotedblleft{}A remembrance; memory\textquotedblright{}.) Say it once more.
\end{tcolorbox}
